% board-warping-lenses.tex — generated by experiments/src/paper_figures/board_warping_offsets.jl; do not edit.
% The three-lens session's board warping in one figure: per board (row) the planar board plus the offset vector that each lens's
% calibration gives its points (in the plane the print error, along the normal the height; drawn 50 times enlarged) as a 3-D surface
% (Sigma 14 mm, Tamron 20 mm, fisheye) and, after the three surfaces, per pair of lenses a map over the page of the length of the difference of
% the two lenses' offset vectors (whole offset, 3-D), as it is: nothing fitted or removed. \input it inside a figure* (it follows \linewidth; pgfplots, compat 1.18).
% The maps, left to right: Sigma 14 mm - Tamron 20 mm, Fisheye - Tamron 20 mm, Sigma 14 mm - Fisheye: all three pairs are shown.
% \rcfigdir is the folder this file lies in, as the host document sees it (default src/fig).
\providecommand{\rcfigdir}{src/fig}%
\begingroup
\ifdefined\rcpanelwd\else\newlength\rcpanelwd\fi
\ifdefined\rcdiffwd\else\newlength\rcdiffwd\fi
% three surfaces of width w, then three maps of width 0.4 w; gaps of 4 pt between surfaces, 10 pt between the two blocks, 8 pt between maps,
% 3 pt free after the last map (a line of its text is wider than the map); the row labels 12 pt
\setlength{\rcpanelwd}{\dimexpr(\linewidth-49pt)*100/420\relax}\setlength{\rcdiffwd}{0.4\rcpanelwd}\def\rcpanelwidth{\rcpanelwd}%
\footnotesize\setlength{\tabcolsep}{0pt}\renewcommand{\arraystretch}{1}%
\newcommand{\rcdiffpanel}[1]{{\def\rcpanelwidth{\rcdiffwd}\input{\rcfigdir/board-warping-lenses/#1}}}%
% a map's header: directly over the map (its last line's baseline 3 pt above it), \tiny; it does not add to the row's height (the header row is empty over the maps)
\newcommand{\rcmaphead}[1]{\llap{\smash{\raisebox{\dimexpr0.5\rcdiffwd+3pt\relax}{\makebox[\rcdiffwd]{\tiny\begin{tabular}[b]{@{}c@{}}#1\end{tabular}}}}}}%
% a map's two lines of numbers: directly under the map, \tiny, centred on it
\newcommand{\rcmapstat}[1]{\llap{\raisebox{\dimexpr-0.5\rcdiffwd-1pt-\height\relax}{\makebox[\rcdiffwd]{\tiny\begin{tabular}[t]{@{}c@{}}#1\end{tabular}}}}}%
\begin{tabular}{@{}p{9pt}@{\hspace{3pt}}c@{\hspace{4pt}}c@{\hspace{4pt}}c@{\hspace{10pt}}c@{\hspace{8pt}}c@{\hspace{8pt}}c@{\hspace{3pt}}}
  & \makebox[0pt]{\scriptsize Sigma 14\,mm lens} & \makebox[0pt]{\scriptsize Tamron 20\,mm lens} & \makebox[0pt]{\scriptsize Fisheye lens} &  &  &  \\\noalign{\vskip 2pt}
  \makebox[9pt][c]{\rotatebox[origin=c]{90}{\scriptsize AprilGrid}} & \input{\rcfigdir/board-warping-lenses/bw_ag_14mm} & \input{\rcfigdir/board-warping-lenses/bw_ag_tamron} & \input{\rcfigdir/board-warping-lenses/bw_ag_fisheye} & \rcdiffpanel{bw_ag_14mm-minus-bw_ag_tamron}\rcmaphead{Sigma 14\,mm $-$\\Tamron 20\,mm}\rcmapstat{RMS 0.53\,mm\\max 0.95\,mm} & \rcdiffpanel{bw_ag_fisheye-minus-bw_ag_tamron}\rcmaphead{Fisheye $-$\\Tamron 20\,mm}\rcmapstat{RMS 0.69\,mm\\max 1.18\,mm} & \rcdiffpanel{bw_ag_14mm-minus-bw_ag_fisheye}\rcmaphead{Sigma 14\,mm $-$\\Fisheye}\rcmapstat{RMS 0.30\,mm\\max 0.97\,mm} \\\noalign{\vskip 5pt}
  \makebox[9pt][c]{\rotatebox[origin=c]{90}{\scriptsize Siemens stars}} & \input{\rcfigdir/board-warping-lenses/bw_si_14mm} & \input{\rcfigdir/board-warping-lenses/bw_si_tamron} & \input{\rcfigdir/board-warping-lenses/bw_si_fisheye} & \rcdiffpanel{bw_si_14mm-minus-bw_si_tamron}\rcmapstat{RMS 0.62\,mm\\max 1.13\,mm} & \rcdiffpanel{bw_si_fisheye-minus-bw_si_tamron}\rcmapstat{RMS 0.43\,mm\\max 0.73\,mm} & \rcdiffpanel{bw_si_14mm-minus-bw_si_fisheye}\rcmapstat{RMS 0.36\,mm\\max 0.79\,mm} \\\noalign{\vskip 6pt}
  \multicolumn{7}{@{}c@{}}{%
    \begin{minipage}[t]{0.49\linewidth}\centering\input{\rcfigdir/board-warping-lenses/key}\\[1pt]{\scriptsize surfaces: the planar board plus the board offsets learned with this lens (in the plane and along the normal), drawn 50 times enlarged; colour: height above the board's plane (positive: towards the camera)}\end{minipage}\hfill
    \begin{minipage}[t]{0.49\linewidth}\centering\input{\rcfigdir/board-warping-lenses/maps-key}\\[1pt]{\scriptsize maps: length of the difference between the board offsets learned with two lenses (in the plane and along the normal); below: RMS and largest value}\end{minipage}} \\
\end{tabular}%
\endgroup
